\begin{definition}[Discrete Function]
A \emph{discrete function} is a function
\[
f:\mathbb{Z} \to \mathbb{R}
\]
(or more generally \(f:\mathbb{N}\to\mathbb{R}\)).

Thus a discrete function assigns a real value to each integer input,
producing a sequence
\[
f(0),\, f(1),\, f(2),\, \dots .
\]
\end{definition}

\begin{remark}
In discrete calculus the independent variable moves in \emph{integer steps}
rather than continuously. Consequently, change is measured between
successive integer inputs. This contrasts with ordinary calculus, where
functions are defined on intervals of real numbers and change is measured
through infinitesimal limits. For this reason, the fundamental operator of
discrete calculus will be the \emph{difference operator}
\[
\Delta f(n)=f(n+1)-f(n),
\]
which plays a role analogous to the derivative in continuous calculus.
\end{remark}

\begin{remark}
A discrete function is simply another way of viewing a \emph{sequence}.
Indeed, a sequence of real numbers
\[
a_0, a_1, a_2, \dots
\]
can be identified with a function whose domain is the integers (or natural
numbers)
\[
f : \mathbb{N} \to \mathbb{R}, \qquad f(n) = a_n.
\]
Thus every sequence determines a discrete function, and every discrete
function determines a sequence. For this reason, both notations are used
interchangeably in discrete mathematics and discrete calculus.

Several equivalent notations commonly appear in the literature:
\[
f(n), \qquad a_n, \qquad (a_n)_{n\in\mathbb{N}}, \qquad \{a_n\}_{n=0}^{\infty}.
\]
All of these describe the same mathematical object: a rule assigning a value
to each integer index. In discrete calculus it is often convenient to use the
functional notation \(f(n)\), since the difference operator will act on the
function much like a derivative acts on functions in continuous calculus.
\end{remark}














